% =====================================================================
%  config/identitas.tex
%  Satu-satunya tempat mengisi metadata skripsi. Semua halaman (sampul,
%  halaman judul, persetujuan, pengesahan, pernyataan, metadata PDF)
%  mengambil nilai dari sini. Ganti teks contoh di dalam kurung kurawal.
% =====================================================================

% --- Judul -----------------------------------------------------------
% Tulis dengan huruf kapital di awal kata; sampul mengubahnya menjadi
% kapital seluruhnya secara otomatis. Jangan memasang \\ untuk memenggal
% baris: pemenggalan dibuat seimbang oleh kelas.
% Istilah asing WAJIB dicetak miring dengan \textit{...} atau \istilah{...},
% dan tetap akan dicetak miring secara otomatis pada sampul luar.
\judul{Evaluasi Kegunaan Aplikasi Layanan Administrasi Desa Menggunakan \textit{System Usability Scale}}
\judulinggris{Usability Evaluation of a Village Administration Service Application Using the \textit{System Usability Scale}}

% --- Mahasiswa -------------------------------------------------------
% Nama mahasiswa ditulis lengkap. Sesuai aturan prodi, penulisan nama mahasiswa
% pada sampul luar dan halaman judul otomatis dicetak menggunakan huruf kapital semua.
\nama{Nama Lengkap Mahasiswa}
\npm{0000000000}
% Label nomor induk yang dicetak di sampul dan lembar lain (NPM atau NIM).
\labelnomorinduk{NPM}

% --- Waktu dan tempat ------------------------------------------------
\tahun{2026}
\tanggalpersetujuan{1 Desember 2026}% tanggal pembimbing menyetujui naskah diuji
\tanggalujian{15 Desember 2026}%      tanggal Sidang Akhir Skripsi
\tanggalpengesahan{29 Desember 2026}% tanggal pengesahan naskah final
% \kota{Sampit}%                      bawaan kelas: Sampit

% --- Dosen Pembimbing (Prodi Sistem Informasi: 1 Orang Pembimbing Tunggal) ---
% Di Program Studi Sistem Informasi UNDA, dosen pembimbing berjumlah 1 orang.
\pembimbing{Nama Dosen Pembimbing, Gelar}
\nuptkpembimbing{0000000000000000}
% Perintah alias untuk kompatibilitas:
\pembimbingutama{Nama Dosen Pembimbing, Gelar}
\nuptkpembimbingutama{0000000000000000}
% Pembimbing pendamping dikosongkan karena standar prodi adalah pembimbing tunggal:
\pembimbingpendamping{}
\nuptkpembimbingpendamping{}

% --- Dewan Penguji Sidang Akhir (2 Orang Penguji) -------------
% Standar naskah skripsi final Prodi SI UNDA: 2 orang penguji
% (Ketua Penguji / Penguji I dan Anggota Penguji / Penguji II).
\ketuapenguji{Nama Ketua Penguji, Gelar}% Ketua Penguji (Penguji I)
\nuptkketuapenguji{0000000000000000}
\anggotapenguji{Nama Anggota Penguji, Gelar}% Anggota Penguji (Penguji II)
\nuptkanggotapenguji{0000000000000000}
% Alias kompatibilitas:
\pengujisatu{Nama Ketua Penguji, Gelar}
\nuptkpengujisatu{0000000000000000}
\pengujidua{Nama Anggota Penguji, Gelar}
\nuptkpengujidua{0000000000000000}

% --- Pimpinan Program Studi (Pengesahan Naskah Skripsi Final) --------
% Sesuai ketentuan prodi, naskah skripsi disahkan oleh Ketua Program Studi.
\kaprodi{Nama Ketua Program Studi, Gelar}
\nuptkkaprodi{0000000000000000}
% Alias kompatibilitas:
\ketuaprogramstudi{Nama Ketua Program Studi, Gelar}
\nuptkketuaprogramstudi{0000000000000000}

% --- Pernyataan Alat Bantu Kecerdasan Buatan (AI Tools) ------------
% Cantumkan seluruh alat bantu AI yang digunakan untuk penyuntingan naskah/debugging:
\aitools{ChatGPT, Claude, Grammarly, GitHub Copilot}

% --- Repositori Kode Sumber Utama dan Mendeley Data -----------------
% Sesuai ketentuan prodi, kode sumber utama wajib dilampirkan tautan GitHub-nya,
% dan dataset/data pendukung diunggah ke Mendeley Data (atau repositori riset resmi):
\tautangithub{https://github.com/mahasiswa/skripsi-sisfo-unda}
\tautanmendeley{https://data.mendeley.com/datasets/sample-dataset/1}

% --- Nilai tetap prodi (sudah diisi kelas, ubah hanya bila berubah) ---
% \programstudi{Sistem Informasi}
% \fakultas{Ilmu Komputer}
% \universitas{Universitas Darwan Ali}
% \gelar{Sarjana Komputer}
